\begin{table}[H]
\centering
\caption{问题二关键约束复算与裕度}
\label{tab:q2-constraint-audit}
\small
\begin{tabular}{lrrrrc}
\toprule
约束族 & 实际值 & 界值 & 单位 & 最小裕度 & 状态 \\
\midrule
最小轮间休息 & 63.000 & 60.000 & 小时 & 5.0\% & 通过 \\
场馆承办下界 & 0.000 & 0.000 & 剩余场 & 0.0\% & 通过 \\
场馆承办上界 & 0.000 & 0.000 & 剩余场 & 0.0\% & 通过 \\
场馆单日本地容量 & 1.000 & 1.000 & 最大利用率 & 0.0\% & 通过 \\
场馆安保等级 & 0.000 & 0.000 & 等级余量 & 0.0\% & 通过 \\
时段转播容量 & 0.000 & 0.000 & 剩余场 & 0.0\% & 通过 \\
第三轮高安保日容量 & 0.000 & 0.000 & 比例余量 & 0.0\% & 通过 \\
黄金时段公平极差 & 1.000 & 2.000 & 场 & 50.0\% & 通过 \\
\bottomrule
\end{tabular}
\vspace{2pt}
\noindent\begin{minipage}{0.96\textwidth}\footnotesize 注：裕度为按各约束自身上限归一化后的最小剩余比例；零表示约束在最紧样本上取等。\end{minipage}
\end{table}
